\begin{graphicspathcontext}{{./chapters/ai/imgs/},{./chapters/ai/imgs/auto/},\old}

\sidecite{Brynjolfsson2014,Acemoglu2018,Freedman2018,Zuboff2019}
\begin{frame}[t]{{Economic impacts:} employment, productivity, inequalities}
	\vspace{-.25cm}
	\begin{columns}[T]
		\begin{column}{.5\linewidth}
			\begin{block}{Opportunities}
				\begin{itemize}
				\item Handling of risky tasks
				\item Decision support
				\item Productivity gains%\cite{Brynjolfsson2014}
				\item New value-added tasks
				\end{itemize}
			\end{block}
		\end{column}
		\begin{column}{.5\linewidth}
			\begin{block}{Risks}
				\begin{description}
				\item[Job losses] repetitive tasks, certain intellectual professions%\cite{Acemoglu2018}
				\item[Concentration of wealth] large companies mastering AI%\cite{Zuboff2019}
				\item[Precarization of workers] algorithmic surveillance, micro-work%\cite{Freedman2018}
				\end{description}
			\end{block}
		\end{column}
	\end{columns}
	\vspace{.5cm}
	\alertbox{\Emph{Central question}: How to share the benefits of AI without widening inequalities?}
	\alertbox*{\emph{Avenues:} continuous training, specific taxation, universal income...}
\end{frame}

\end{graphicspathcontext}

\endinput
